% ============================================================================
% RELATED WORK (REVISED v2 - With RouteLLM Calibration Bottleneck)
% ============================================================================

\section{Related Work}
\label{sec:related}

Our work builds upon and learns from three research areas: adaptive model selection, cost-aware inference optimization, and transfer learning for decision-making systems. We position \ours{} as a complementary alternative that prioritizes accessibility and ease of maintenance in rapidly evolving model ecosystems.

\subsection{LLM Routing and Model Selection}

\paragraph{Cascading Systems.} FrugalGPT~\cite{chen2023frugalgpt} pioneered cost-optimized LLM routing through sequential model invocation with learned verification thresholds, demonstrating impressive cost-quality trade-offs on fixed model pools. The system validates that \emph{adaptive model selection works}---a foundational insight upon which we build. However, operational requirements limit accessibility: users often need hundreds to thousands of labeled examples to calibrate chains and train domain-specific scoring functions (typically DistilBERT regressors)~\cite{chen2023frugalgpt}. For users without annotated datasets or ML expertise to design verifiers, this setup cost is prohibitive.

We learn from FrugalGPT's strength in cascading reliability, incorporating selective verification in our Hybrid Mode (\S\ref{sec:hybrid}). However, we address the calibration barrier through \emph{shippable priors}---pre-trained covariance matrices that enable zero-calibration deployment. This reduces setup from days (data collection + scorer training) to minutes (prior download), expanding accessibility from ML teams to general programmers.

\paragraph{Static Classification Systems.} RouteLLM~\cite{ong2024routellm} demonstrates that preference-based learning effectively captures user intent for routing between two models. By training BERT classifiers on the RouterBench dataset from LMSYS Chatbot Arena, the system achieves strong performance for general conversational tasks. This validates supervised learning as a viable routing paradigm.

However, RouteLLM faces a \emph{calibration bottleneck} that limits scalability. As a supervised learning approach, adding new models requires collecting new preference pairs at scale (typically in the thousands) and retraining the classifier to incorporate the new model's characteristics~\cite{ong2024routellm}. This constitutes an $O(N)$ maintenance cost per model addition, where $N$ is the dataset size. For organizations without ML infrastructure, the expertise required to curate domain-specific datasets, evaluate outputs, and retrain classifiers creates an operational barrier. In markets where 10+ models launch monthly, this recalibration treadmill becomes unsustainable: by the time engineers update the router, several new alternatives have released.

Our reinforcement learning approach addresses this through $O(1)$ model registration: new models are added via config updates (seconds) and evaluated autonomously through online exploration. This decouples maintenance cost from pool size, enabling users to track market evolution without ML expertise or dedicated calibration pipelines.

\paragraph{Complementary Strengths.} We emphasize that both approaches have merit:
\begin{itemize}[leftmargin=*,nosep]
    \item \textbf{Static systems excel when:} Model pools are stable (2--3 fixed options), users have labeled datasets, and ML teams can perform periodic recalibration.
    \item \textbf{Adaptive systems excel when:} Model ecosystems evolve rapidly (80+ options, weekly releases), users lack calibration data, or maintenance must be sustainable without ML specialists.
\end{itemize}

\begin{table}[t]
\centering
\footnotesize
\caption{\textbf{Maintenance Requirements Comparison.} We address the calibration bottleneck through online learning.}
\label{tab:maintenance_comparison}
\begin{tabular}{lccc}
\toprule
\textbf{Operation} & \textbf{FrugalGPT} & \textbf{RouteLLM} & \textbf{Ours} \\
\midrule
Add New Model & Re-run benchmarks & Retrain classifier & Register + explore \\
Time Required & 1--3 days & 1--3 days & 5 minutes \\
Cost per Model & \$50--200 & \$50--200 & \$0 (online) \\
Expertise & High & Medium--High & None \\
Scaling & $O(N)$ & $O(N)$ & $O(1)$ \\
\bottomrule
\end{tabular}
\end{table}

\paragraph{A Critical Limitation: Confident Failure.} Through systematic evaluation (\S\ref{sec:confident_failure}), we identify a failure mode in reactive verification systems: scoring functions incorrectly validate 35\% of erroneous outputs from cost-optimized models when outputs are plausible but semantically incorrect. Our predictive architecture mitigates this by routing high-uncertainty prompts directly to frontier models before generation, preemptively avoiding the pathway where confident failures occur. This insight suggests that ex-ante risk assessment complements ex-post verification.

\subsection{Multi-Armed Bandits and Online Learning}

\paragraph{Contextual Bandits.} Contextual multi-armed bandits~\cite{li2010contextual,chu2011contextual} provide a principled framework for sequential decision-making under uncertainty, with LinUCB~\cite{li2010contextual} achieving provably optimal regret bounds under linear reward assumptions. However, standard bandits suffer from cold-start inefficiency: convergence requires hundreds to thousands of samples, translating to substantial exploration costs (\$50--200) that price out resource-constrained users.

\paragraph{Our Contribution: Warm-Start Mechanism.} We address cold-start inefficiency through offline expert distillation (\S\ref{sec:distillation}). By pre-training covariance matrices on teacher supervision, we reduce day-one regret by 64.6\% (\S\ref{sec:rq1}) while preserving the bandit's adaptive plasticity. This differs from standard transfer learning~\cite{taylor2009transfer} by explicitly encoding uncertainty structure (covariance) rather than point estimates alone, enabling informed exploration from the first query without user-provided calibration data.

\paragraph{Plasticity Under Drift.} A key advantage of online learning over static systems is autonomous adaptation to concept drift. We validate this through poisoned prior experiments (\S\ref{sec:rq2}), demonstrating that \ours{} successfully corrects teacher bias and model capability changes without manual intervention. In contrast, static classifiers (RouteLLM) cannot detect when provider updates degrade model performance; they continue routing to outdated preferences until manually retrained. This self-correction capability is critical for production systems where model capabilities evolve continuously.

\subsection{LLM Benchmarking and Evaluation}

\paragraph{Static Benchmarks.} Comprehensive evaluation frameworks like HELM~\cite{liang2022holistic}, BIG-bench~\cite{srivastava2022beyond}, and LMSYS Chatbot Arena~\cite{zheng2023judging} provide valuable model comparisons. We leverage these benchmarks for offline prior distillation and external validation (\S\ref{sec:rq1}).

\paragraph{The Benchmark Trap.} However, initializing routing policies from public benchmark leaderboards yields worse performance than random initialization (Table~\ref{tab:benchmark_trap}). This occurs because benchmark-optimal models diverge systematically from deployment-optimal models: benchmarks evaluate broad capability across standardized distributions, whereas effective routing requires utility optimization under deployment-specific constraints (cost, latency, local task distributions). This finding suggests that routing systems require contextual adaptation rather than global rankings.

\subsection{Design Philosophy: Complementary Alternatives}

Rather than positioning \ours{} as a replacement for existing systems, we view it as a complementary alternative optimized for different operational contexts:

\begin{itemize}[leftmargin=*,nosep]
    \item \textbf{FrugalGPT excels when:} Users have labeled calibration data, domain expertise to design scorers, stable prompt distributions where upfront optimization costs amortize over long deployments, and ML teams to perform periodic maintenance.
    \item \textbf{RouteLLM excels when:} Users have access to labeled preference data, domains align with RouterBench (general chat), and model pools are small and stable (2--3 options).
    \item \textbf{BanditGPT excels when:} Users lack calibration data or ML expertise, face rapidly evolving model ecosystems (80+ options, weekly releases), require sustainable maintenance without dedicated ML infrastructure, or prioritize immediate deployment over asymptotic optimality.
\end{itemize}

Our core insight is that \emph{accessibility requires ease of maintenance, not just ease of deployment}. A system that deploys quickly but becomes outdated within months fails to democratize access. By eliminating recalibration requirements through online learning, we enable sustainable routing for users without ML teams---expanding adaptive model selection from specialized deployments to mainstream adoption.

\subsection{Positioning Summary}

\begin{table}[t]
\centering
\footnotesize
\caption{\textbf{Learning from Prior Work.} \ours{} combines strengths while addressing accessibility barriers.}
\label{tab:related_comparison}
\begin{tabular}{p{2.0cm}p{2.6cm}p{2.2cm}}
\toprule
\textbf{System} & \textbf{Strength (We Learn)} & \textbf{Limitation (We Address)} \\
\midrule
FrugalGPT & High reliability via cascading & Heavy setup; $O(N)$ maintenance \\
RouteLLM & Preference learning & Static; $O(N)$ recalibration \\
Std.\ Bandits & Principled exploration & Cold-start (\$50--200) \\
Leaderboards & Global capability rankings & Cost-agnostic; not contextual \\
\bottomrule
\end{tabular}
\end{table}

The fundamental innovation enabling \ours{}'s accessibility is the combination of two architectural decisions: (1)~\emph{shippable priors} that eliminate calibration requirements for deployment, and (2)~\emph{online exploration} that eliminates recalibration requirements for maintenance. This dual mechanism addresses both the setup barrier (users without datasets) and the sustainability barrier (users without ML teams). By reducing economic barriers (61\% cost reduction), expertise barriers (zero calibration data), and maintenance barriers ($O(1)$ vs.\ $O(N)$ scaling), we aim to expand adaptive routing from specialized ML teams to the broader community of students, researchers, and practitioners.

